\documentclass[10pt,a4paper]{amsart}
\usepackage[margin=19mm]{geometry}
\usepackage{amsmath,amssymb,mathtools}
\usepackage[scr=rsfs,cal=euler]{mathalfa}
\usepackage{xfrac}
\usepackage[unicode,hidelinks,bookmarks=false]{hyperref}
\newcommand{\ie}{i.e.\ }
\newcommand{\cf}{cf.\ }
\newcommand{\resp}{respectively\ }
\newcommand{\Subsection}{\subsection}
\providecommand{\nobreakdash}{\nobreak\hbox{-}}
\setlength{\emergencystretch}{2em}
\multlinegap0pt
\usepackage{mathtools}
\usepackage{enumitem}
\hypersetup{colorlinks=true,citecolor=blue,linkcolor=blue,urlcolor=blue}
\setlength{\emergencystretch}{2em}
\newtheorem{theorem}{Theorem}[section]
\newtheorem{lemma}[theorem]{Lemma}
\newtheorem{proposition}[theorem]{Proposition}
\newtheorem{corollary}[theorem]{Corollary}
\newtheorem{problem}[theorem]{Problem}
\newtheorem{conjecture}[theorem]{Conjecture}
\newtheorem{definition}[theorem]{Definition}
\newtheorem{example}[theorem]{Example}
\newtheorem{remark}[theorem]{Remark}
\newcommand{\R}{\mathbb R}
\newcommand{\eps}{\varepsilon}
\newcommand{\W}{\mathcal W}
\newcommand{\dd}{\,d}
\newcommand{\norm}[1]{\left\lVert #1\right\rVert}
\newcommand{\diag}{\operatorname{diag}}
\newcommand{\SO}{\operatorname{SO}}
\begin{document}
\title[Frenet turns]{Frenet turns}
\author{Boris Shapiro}
\date{}
\maketitle
\noindent\emph{Source and licence.} Frenet turns. Version arXiv:2606.10606v1 (CC BY 4.0). This material is distributed under the Creative Commons Attribution 4.0 International licence, http://creativecommons.org/licenses/by/4.0/, as recorded in the arXiv OAI licence metadata for this exact version and shown on the abstract page https://arxiv.org/abs/2606.10606v1. Full rights notice: licenses/arxiv-2606.10606-CC-BY-4.0.txt.\par\smallskip
\noindent\textit{Editorial scope.} Counted unit: the original \texttt{theorem} environment \texttt{thm:even-constant-openings} at source lines 570--589 together with its complete proof at lines 591--618; the delivered document keeps source lines 340--619 so that every referenced label and every citation resolves inside it. Native editable source is the paper's own concatenated TeX; the AMS wrapper is editorial formatting only. No publisher class, upstream script or unrelated artwork is redistributed.
\begin{lemma}[spherical Fenchel obstruction]\label{lem:spherical-obstruction}
Let $v:S^1\to S^{N-1}$ be a closed rectifiable curve and let $a:S^1\to(0,\infty)$
be integrable.  If
\[
        \int_{S^1}a(t)v(t)\,dt=0,
\]
then the spherical length of $v$ is at least $2\pi$.
\end{lemma}

\begin{proof}
The condition says that the origin lies in the convex hull of the image of $v$.
Therefore the image is not contained in any open hemisphere.  The spherical form of
Fenchel's theorem gives length at least $2\pi$; see, for example,
\cite{Fenchel,Milnor}.
\end{proof}

\begin{theorem}[control obstruction]\label{thm:strong-control-obstruction}
Let $n\ge4$, and let $E=(e_1,\ldots,e_n)$ be a closed positive Frenet frame in
$\R^n$.  Then
\begin{equation}\label{eq:last-bound}
        \int_0^T\sqrt{u_{n-2}(t)^2+u_{n-1}(t)^2}\,dt\ge2\pi .
\end{equation}
Consequently no fixed normally framed multiply traversed plane curve can be
approached by positive Frenet controls in any norm topology implying
\[
        \|u_{n-2}\|_{L^1}+\|u_{n-1}\|_{L^1}\longrightarrow0.
\]
\end{theorem}

\begin{proof}
The last Frenet equation is
\[
        e_n'=-u_{n-1}e_{n-1}.
\]
Since $e_n(T)=e_n(0)$,
\[
        \int_0^T u_{n-1}(t)e_{n-1}(t)\,dt=0.
\]
Apply Lemma~\ref{lem:spherical-obstruction} to $v=e_{n-1}$ and $a=u_{n-1}$.  Since
\[
        e_{n-1}'=-u_{n-2}e_{n-2}+u_{n-1}e_n,
\]
the spherical length of $e_{n-1}$ is exactly the left-hand side of
\eqref{eq:last-bound}.  This proves the inequality.  The final claim follows because
the boundary controls of a normally framed plane curve have $u_{n-2}=u_{n-1}=0$.
\end{proof}

\begin{remark}\label{rem:why-not-number}
Theorem~\ref{thm:strong-control-obstruction} is an obstruction, not a useful number
of turns.  With a fixed normal frame in dimensions $n\ge4$, the strong control answer
is simply that no finite number of turns suffices.  A meaningful version of
Agrachev's problem should retain a turn-counting datum while avoiding the artificial
requirement that the limiting normal frame be fixed.
\end{remark}

\section{Decorated turn vectors in dimension four}

We now describe such a turn-counting replacement in the first higher-dimensional
case.  Work on $0\le t\le2\pi$ and identify the initial Frenet frame with the identity.
In dimension four the degenerate boundary controls
\begin{equation}\label{eq:decorated-control}
        (u_1,u_2,u_3)=(p,0,q),
        \qquad p,q\in\mathbb Z_{\ge0},\quad p\ge1,
\end{equation}
produce the boundary frame
\[
        E_{p,q}(t)=R_{12}(pt)\oplus R_{34}(qt).
\]
The base curve has $p$ tangent turns, while the normal $2$-frame has $q$ turns.  We
call $(p,q)$ the decorated turn vector.

\begin{definition}\label{def:accessible-turn-vector}
A turn vector $(p,q)$ is called \emph{strongly accessible} if, for every $\eta>0$,
there exist positive controls $u_1,u_2,u_3$ on $[0,2\pi]$ satisfying the closed-frame
and closed-curve endpoint conditions
\begin{equation}\label{eq:endpoint-conditions}
        E(2\pi)=I,
        \qquad
        \int_0^{2\pi}E(t)e_1\,dt=0,
\end{equation}
and
\[
        \|u_1-p\|_{L^\infty}+\|u_2\|_{L^\infty}+\|u_3-q\|_{L^\infty}<\eta .
\]
\end{definition}

The fixed-normal-frame case is $(p,0)$.  Theorem~\ref{thm:strong-control-obstruction}
shows that no $(p,0)$ with $p\ge1$ is strongly accessible.  Thus the normal turn is
not a decoration in name only; it is forced by the last-vector obstruction.  The next
result gives the opposite direction for all nonresonant constant turn vectors.

\begin{theorem}[constant-control opening in dimension four]\label{thm:constant-control-opening}
Every turn vector $(p,q)$ with $p,q\in\mathbb Z_{>0}$ and $p\ne q$ is strongly
accessible.  More precisely, for all sufficiently small $\eps>0$ there exist positive
constants $a_\eps,c_\eps$ such that the constant controls
\[
        (u_1,u_2,u_3)=(a_\eps,\eps,c_\eps)
\]
satisfy \eqref{eq:endpoint-conditions}, and
\[
        a_\eps\to p,
        \qquad
        c_\eps\to q
        \qquad(\eps\to0).
\]
\end{theorem}

\begin{proof}
For constant controls $(a,b,c)$, write the Frenet equations as $E'=EB$, where
\[
        B=\begin{pmatrix}
        0&-a&0&0\\
        a&0&-b&0\\
        0&b&0&-c\\
        0&0&c&0
        \end{pmatrix}.
\]
The eigenvalues of $B$ are $\pm i\lambda_1,\pm i\lambda_2$, where
\begin{equation}\label{eq:frequency-relations}
        \lambda_1^2+\lambda_2^2=a^2+b^2+c^2,
        \qquad
        \lambda_1^2\lambda_2^2=a^2c^2 .
\end{equation}
We prescribe $\lambda_1=p$, $\lambda_2=q$, and set $b=\eps$.  Then $a,c$ must
satisfy
\[
        a^2+c^2=p^2+q^2-\eps^2,
        \qquad
        ac=pq .
\]
Equivalently, $a^2$ and $c^2$ are the roots of
\[
        z^2-(p^2+q^2-\eps^2)z+p^2q^2=0 .
\]
Since $p\ne q$, the discriminant is positive for all sufficiently small $\eps$, and
the two roots tend to $p^2$ and $q^2$.  This gives positive $a_\eps,c_\eps$ with the
required limits.

By the skew-symmetric spectral theorem, the frame is closed because both frequencies
are integers; in other words $\exp(2\pi B)=I$.  Since $a_\eps c_\eps>0$, the matrix
$B$ is invertible, and therefore
\[
        \int_0^{2\pi}\exp(tB)e_1\,dt
        =B^{-1}(\exp(2\pi B)-I)e_1=0.
\]
Thus the base curve is closed as well.
\end{proof}

\begin{corollary}\label{cor:first-one-turn}
The decorated turn vector $(1,2)$ is strongly accessible.  It is the smallest
nonresonant decorated vector with one tangent turn.
\end{corollary}

\begin{remark}[the resonant case]\label{rem:resonance}
The proof of Theorem~\ref{thm:constant-control-opening} fails when $p=q$.  In fact,
there is no constant-control opening with the same double frequency, because a
positive middle control splits the double eigenvalue.  This does not rule out
strong accessibility by time-dependent controls.  The resonant vectors $(p,p)$ are
therefore small concrete instances of Agrachev's one-sided endpoint problem.
\end{remark}

\section{A revised Agrachev problem}

The preceding discussion suggests replacing the fixed-normal-frame question by the
following decorated turn problem.  This is not a reformulation of Agrachev's original
number $\mu(n)$; rather, it is a nearby boundary-value problem designed to keep the
turn-counting feature while removing the artificial obstruction caused by freezing
the normal frame.

\begin{problem}[decorated Agrachev problem in dimension four]\label{prob:decorated-R4}
Determine all pairs $(p,q)\in\mathbb Z_{\ge0}^2$, $p\ge1$, for which the degenerate
Frenet trajectory
\[
        R_{12}(pt)\oplus R_{34}(qt),
        \qquad 0\le t\le2\pi,
\]
is in the strong control closure of positive closed Frenet trajectories satisfying
\eqref{eq:endpoint-conditions}.  Equivalently, determine which degenerate controls
$(p,0,q)$ can be opened to positive controls $u_1,u_2,u_3>0$.
\end{problem}

The present note proves the two first facts about this problem:
\begin{enumerate}[label=\textup{(\roman*)}]
\item the undecorated vectors $(p,0)$ are not strongly accessible;
\item all nonresonant vectors $(p,q)$ with $p,q>0$ and $p\ne q$ are strongly
accessible by constant controls.
\end{enumerate}
Thus the problem has genuine turn-counting content in both directions.  It also
separates the topological and variational parts of Agrachev's original question.  The
number of tangent turns alone is not sufficient in dimension four; the normal turn
must be part of the boundary datum.

The dimension-four statement is the first member of a higher-dimensional pattern.
We formulate it next.

\begin{definition}
For $n=2r$, a \emph{decorated turn vector} is a vector
\[
        \mathbf p=(p_1,\ldots,p_r)\in\mathbb Z_{>0}^r .
\]
It represents the degenerate boundary frame
\begin{equation}\label{eq:even-decorated-frame}
        E_{\mathbf p}(t)=R_{12}(p_1t)\oplus R_{34}(p_2t)\oplus\cdots\oplus
        R_{2r-1,2r}(p_rt),\qquad 0\le t\le2\pi .
\end{equation}
Equivalently, in the constant Frenet matrix the odd controls are
\[
        u_{2j-1}=p_j,
        \qquad j=1,\ldots,r,
\]
while the even controls $u_2,u_4,\ldots,u_{2r-2}$ vanish.  We say that $\mathbf p$
is \emph{strongly accessible} if this degenerate datum is in the $L^\infty$-closure
of positive closed Frenet controls with closed base curve and closed frame.
\end{definition}

Here and below the tridiagonal Frenet matrix associated with constant controls
$a_1,\ldots,a_{n-1}$ is
\begin{equation}\label{eq:general-tridiagonal-B}
        B(a_1,\ldots,a_{n-1})=
        \begin{pmatrix}
        0&-a_1&0&\cdots&0\\
        a_1&0&-a_2&\cdots&0\\
        0&a_2&0&\cdots&0\\
        \vdots&\vdots&\vdots&\ddots&-a_{n-1}\\
        0&0&0&a_{n-1}&0
        \end{pmatrix}.
\end{equation}
The dimension-four theorem above is the case $r=2$ of the following elementary
inverse-spectral observation.


\begin{theorem}[even-dimensional constant openings]\label{thm:even-constant-openings}
Let $n=2r$.  If
\[
        \mathbf p=(p_1,\ldots,p_r)\in\mathbb Z_{>0}^r
\]
has pairwise distinct entries, then $\mathbf p$ is strongly accessible.  More
precisely, for every sufficiently small choice of positive numbers
\[
        b_1,\ldots,b_{r-1},\qquad b_j>0,
\]
there exist positive numbers $a_1,\ldots,a_r$ with $a_j\to p_j$ as
$b_1,\ldots,b_{r-1}\to0$ such that the constant controls
\begin{equation}\label{eq:even-opening-controls}
        u_{2j-1}=a_j,
        \qquad
        u_{2j}=b_j
        \quad (j=1,\ldots,r-1),
\end{equation}
generate a closed positive Frenet curve and a closed Frenet frame on $[0,2\pi]$.
\end{theorem}

\begin{proof}
At $b_1=\cdots=b_{r-1}=0$, the matrix is block diagonal, with $2\times2$ blocks
rotating with frequencies $a_1,\ldots,a_r$.  Since the limiting frequencies
$p_1,\ldots,p_r$ are distinct, the positive eigenfrequencies of the skew-symmetric
matrix are smooth functions
\[
        \lambda_j=\lambda_j(a_1,\ldots,a_r,b_1,\ldots,b_{r-1})
\]
near the boundary point, after ordering them near $p_j$.  At the boundary,
\[
        \lambda_j(a_1,\ldots,a_r,0,\ldots,0)=a_j,
\]
so the Jacobian of $(\lambda_1,\ldots,\lambda_r)$ with respect to
$(a_1,\ldots,a_r)$ is the identity at $(p_1,\ldots,p_r,0,\ldots,0)$.  The implicit
function theorem therefore gives unique smooth positive functions $a_j=a_j(b)$,
near $p_j$, such that
\[
        \lambda_j(a(b),b)=p_j,
        \qquad j=1,\ldots,r.
\]
For these controls all frequencies of $B$ are integers; hence $\exp(2\pi B)=I$.
Because $n$ is even and all frequencies are nonzero, $B$ is invertible.  Therefore
\[
        \int_0^{2\pi}\exp(tB)e_1\,dt
        =B^{-1}(\exp(2\pi B)-I)e_1=0.
\]
Thus the Frenet frame and the base curve are both closed.
\end{proof}


\begin{thebibliography}{99}
\bibitem[Fenchel]{Fenchel} W.~Fenchel, \emph{On the differential geometry of closed space curves}, Bull. Amer. Math. Soc. 57 (1951), 44--54.
\bibitem[Milnor]{Milnor} J.~Milnor, \emph{On total curvatures of closed space curves}, Math. Scand. 1 (1953), 289--296.
\end{thebibliography}
\end{document}
